\paragraph{What the routing changes mean.}
The pre-specified primary test confirms on the held-out split: the medical model's
median question dominance ratio is lower ($5.36 \to 4.15$, $p{=}0.006$; the mean does not
move in that direction, see Section~\ref{sec:confirmation}), so relative to its
language prior it leans less on the question text when it answers. The effect is
\emph{small} (Cliff's $\delta{=}0.18$): reliable on fresh samples, but modest in
size. It is also much smaller than the drop seen in an early 50-sample pilot that did not
hold up; we frame it as a consistent small shift, not a dramatic rerouting. Alongside this, the
confirmed candidates describe \emph{how} the decoder's behaviour changes rather than
simply that it ``uses the image more'': the medical model commits to its first answer
token with much lower entropy, produces smaller closed yes/no logit margins, and
generates lower-probability (more hedged) answer prose; it also shows lower absolute information flow (IF AUC), though that effect survives
only at the raw level after correction. Normalised by each model's answer probability,
MedGemma's flow is about 0.6 times Gemma's in both pathways, so the reduction is general
rather than specific to the image. The pathway-specific evidence is the question
dominance ratio. Notably, the half-layer ``image enters later'' shift
that looked robust on discovery (entry layer $11.0 \to 11.5$, $p_\text{Holm}{=}0.004$)
\emph{did not replicate} on the confirmation split (median 11.0 for both, $p{=}0.43$);
the pre-specified protocol is what lets us retract it cleanly rather than carry it as a finding.

\paragraph{The central open problem: encoder vs.\ decoder.}
MedGemma was adapted with a medical SigLIP vision encoder, and we confirm that the
two encoders differ substantially (layer-0 image-feature linear CKA 0.155; mean
cosine 0.33). The routing changes we measure in the decoder are therefore
\emph{confounded}: they may reflect the decoder reading a different image
representation, a genuinely re-routed decoder, or both. We attempted a correlational
decomposition (regressing the per-sample grounding shift on per-sample encoder
divergence), but it is uninformative here, because the per-sample encoder
divergence is near-constant ($\approx 0.68$ for every sample), leaving no
encoder-similar subset to contrast. We therefore make no causal attribution. The
clean test is an \textbf{encoder swap} (grafting the base encoder onto the medical
decoder and vice versa), which we pre-specify and pursue as dedicated follow-up work.

\paragraph{Scope of the accuracy claim.}
Accuracy is reported for closed-ended yes/no questions only, using the logit verdict
P(yes) vs.\ P(no) at the first answer token. Open-ended questions are not evaluated:
both models generate explanatory prose and, at \texttt{max\_new\_tokens=64}, the
answer often does not appear before the cutoff. The closed-accuracy gap is
underpowered on discovery (59 questions, $+18.6$ points, Fisher $p{=}0.050$ raw) but
large and significant on the held-out split (77.2\% vs.\ 44.3\% on the forced-choice verdict,
$p{=}3.9{\times}10^{-5}$). We are deliberately cautious about its size: the base model
answers near a trivial baseline (at the always-``yes'' rate and below the majority
class), so the model-vs-model difference overstates the medical model's true advantage.
The claim we stand behind is that the medical model clears the majority-class baseline
and the base model does not. Routing metrics do not depend on the answer text and are
reported on all samples.

\paragraph{Discovery vs.\ confirmation.}
The two-split protocol earned its keep. The pre-specified primary test (QDR) was only
borderline on discovery yet confirmed on the held-out split, while the entry-layer
shift showed the opposite pattern: robust on discovery, absent on confirmation. Had
we reported the discovery split alone, we would have published the entry-layer effect
and hedged on the routing claim; the confirmation split reverses both calls. Reporting
candidates and nulls alike, with raw and Holm-adjusted $p$-values, is a deliberate
guard against the selective reporting that a small single-split study invites.

\paragraph{Computational cost.}
Each sample requires a full layer sweep per block type on top of a baseline pass. On
Kaggle T4 GPUs (two per kernel, the free tier), a 100-sample discovery run across two block
types took 5.1 h (Gemma) and 5.6 h (MedGemma) of kernel wall time, and each 150-sample
confirmation run took 7.6 h, split over two GPUs. The method is suited to offline interpretability
analysis rather than online inference.


\paragraph{Limitations.}
This is a single-author study on one benchmark with 250 samples: 100 for discovery and
150 for confirmation. That leaves limited power for small effects, and the primary effect
is small (Cliff's $\delta{=}0.18$). The protocol was written after the discovery analysis
(Section~\ref{sec:hypotheses}), so only the confirmation split is a clean test of the
pre-specified claims. Accuracy is scored on closed yes/no questions only, using logit
verdicts. Open-ended generations are capped at 64 tokens, so they are not scored, and the
accuracy numbers would be unreliable for longer answers. Raising the cap from 8 to 64 tokens
did not produce scorable open-ended answers in our runs, so open-ended accuracy is reported
as zero-coverage rather than as a result. The runs used free Kaggle GPUs (two T4s per run),
which set the sample count and the generation length we could afford. Longer generation and
a full open-ended evaluation are left for follow-up work with more compute. Some Weights \&
Biases logs were not recorded in the cloud, so the Kaggle logs and exported files are the
complete record. The question-dominance ratio is
relative to the model's own language prior, so it does not measure absolute reliance on
the text. The encoder-versus-decoder question is unresolved: the encoder change is
documented, but the correlational analyses cannot attribute the routing change to either
part. We used one model pair, so the findings may not transfer to other medical
fine-tunes or other benchmarks.
